\subsection{Leaf recurrence}
%
Let $i\in V(Q)$ be a leaf vertex of a quiver $Q$, connected to some other vertex $j\in V(Q)$. We call $i$ a \emph{\Louiseleaf} if both quivers $Q-\{i\}$ and $Q-\{i,j\}$ are \Louise.  This automatically implies that $Q$ itself is \Louise.

%

\begin{proposition}\label{prop:Louise_leaf}
Fix a field $\Ga$, and consider cluster varieties over $\Ga$.
Let $\Qice$ be a really full rank ice quiver with $m$ frozen vertices and mutable part $Q$. Let $i\in V(Q)$ be a \Louiseleaf in $Q$ connected to $j\in V(Q)$. Then the subvariety of $\XQice$ given by $x_i=0$ is isomorphic to $\Ga\times \Xcal(\Qice')$, where $\Qice'$ is a really full rank quiver with $m$ frozen vertices and mutable part $Q':=Q-\{i,j\}$.
\end{proposition}
\begin{proof}%
Our goal is to construct a quiver $\Qice'$ satisfying the above assumptions and an isomorphism
\begin{equation}\label{eq:leaf1}
  \{\bx\in\Xcal(\Qice)\mid x_i=0\} \cong \F\times \Xcal(Q').
\end{equation}

  By Proposition~\ref{prop:Mulcover}, the subvariety of $\Xcal(\Qice)$ given by $x_i=0$ is contained inside the subvariety of $\XQice$ given by $x_j\neq0$. The quiver $Q-\{j\}$ is the disjoint union of $Q'$ and a single isolated vertex. By assumption, $Q'$ is \Louise and thus so is $Q-\{j\}$. By Proposition~\ref{prop:Mulloc}, we have $\AQice[x_j^{-1}]\cong \AX(\Qicefr[j])$. We therefore find
\begin{equation}\label{eq:leaf2}
    \{\bx\in\Xcal(\Qice)\mid x_i=0\} \cong \{\bx\in\Xcal(\Qicefr[j])\mid x_i=0\}.
\end{equation}

%
Let $A= \AX(\Qicefr[j] - \{i\})[x_i^{\pm 1}]$ be obtained from the cluster algebra $ \AX(\Qicefr[j] - \{i\})$ by adjoining $x_i$ and its inverse.  The cluster algebra $\AX(\Qicefr[j])$ is isomorphic 
%
 to the subalgebra of $A$ generated by $\AX(\Qicefr[j] - \{i\})$, the cluster variable $x_i$ and the mutation $x'_i$, which is of the form $f/x_i$ where $f \in \AX(\Qicefr[j] - \{i\})$.  

Let $\Xcal':=\Xcal(\Qicefr[j] - \{i\})$.  It follows that 
\begin{equation}\label{eq:leaf2.5}
  \Xcal(\Qicefr[j])\cong \{\left((x_i,x_i'),\bx'\right)\in\Ga^2\times\Xcal'\mid x_ix_i'=M+x_jM'\},
\end{equation}
where $M,M'$ are monomials in the frozen variables of $\Qicefr[j] - \{i\}$. We therefore see that the right-hand side of~\eqref{eq:leaf2} is given by 
\begin{equation}\label{eq:leaf3}
  \{\bx\in\Xcal(\Qicefr[j])\mid x_i=0\}
\cong \{\left((x_i,x_i'),\bx'\right)\in\Ga^2\times\Xcal'\mid x_i=0 \text{ and } x_ix_i'=M+x_jM'\}.
\end{equation}
We see that the variable $x_i'$ is free, and thus the right-hand side of~\eqref{eq:leaf3} is isomorphic to the direct product of $\Ga$ and 
\begin{equation} \label{eq:yM'/M=1} 
\left\{\bx'\in \Xcal'\ \middle|\  \frac{x_jM'}{M}=-1\right\}.
\end{equation}
%
%
%
%
%
%
%
%
%
%
%
%
%
%
%
  It remains to show that the locus~\eqref{eq:yM'/M=1} is isomorphic to $\Xcal(\Qice')$ for some quiver $\Qice'$ satisfying the conditions in \cref{prop:Louise_leaf}. By assumption, $\Bice:=\BQice$ is really full rank. Thus, there exist integers $(\alpha_k)_{k\in[n+m]}$ such that 
  \begin{equation*}%
    \sum_{k} \alpha_k \bice_k=e_i,
  \end{equation*}
where $\bice_k$ denotes the $k$-th row of $\Bice$ and $e_i=(0,\dots,0,1,0,\dots,0)$ is the $i$-th standard basis vector in $\Z^n$. 

%
Since $i$ is a leaf, the row $\bice_i$ has a single nonzero entry in the $j$-th column. Let $\Bice\pA$ be obtained from $\Bice$ by removing the $i$-th row and the $j$-th column.\footnote{When removing rows and columns from matrices, we preserve the labels of the remaining rows and columns. For instance, removing row $1$ from $\Bice$ produces a matrix with rows labeled $2,\dots,n+m$.} Then we have
  \begin{equation*}%
    \sum_{k\neq i} \alpha_k \bice\pA_k=\bar e_i,
  \end{equation*}
where $\bar e_i$ is the $i$-th standard basis vector in $\Z^{n-1}$. Moreover, the matrix $\Bice\pA$ is still really full rank.
%
%
%
%
%

Let $\Vfro:=[n+1,n+m]$ be the set of frozen vertices of $\Qice$. Then the set of frozen vertices of $\Qicefr[j]-\{i\}$ is $\Vfro\sqcup\{j\}$.  We have $\gcd(\alpha_k\mid k\in \Vfro\sqcup\{j\})=1$
%
 since all the rows with a nonzero entry in column $i$ of $\Bice$ belong to $\Vfro\sqcup\{j\}$.  By \cite[Section 5]{LS}, replacing a frozen row by its negative, or adding a frozen row to another frozen row produces an isomorphic cluster variety.  After applying a series of such transformations, we obtain a really full rank $(n+m-1)\times(n-1)$ exchange matrix $\Bice\pB$ and a family of integers $(\alpha\pB_k)_{k\neq i}$ satisfying 
   \begin{equation}\label{eq:alpha_kBice''_k=e_i}
    \sum_{k\neq i} \alpha\pB_k \bice\pB_k=\bar e_i,
  \end{equation}
such that moreover for some $\ell\in\Vfro\sqcup\{j\}$ we have $\alpha\pB_\ell=1$. 

Now, let $\Bice\pC$ be the $(n+m-1)\times(n-2)$ matrix obtained from $\Bice\pB$ by further removing the $i$-th column. Thus, $\AX(\Bice\pC)\cong \AX(\Qicefr[j] - \{i\})$, and by construction, $\sum_{k\neq i} \alpha\pB_k\bice\pC_k=0$ is the zero vector in $\Z^{n-2}$.  Thus, by~\cite[Proposition 5.1]{LS}, the numbers $(\alpha\pB_k)_{k\neq i}$ describe a one-parameter subgroup $\Gm$ of the cluster automorphism torus $T=\Aut(\AX(\Bice\pC))$ (see \cite[Section 5]{LS}), with $z\in\Gm$ acting by $x_k\mapsto z^{\alpha\pB_k}x_k$. Equation~\eqref{eq:alpha_kBice''_k=e_i} implies that $z$ acts on the monomial $\frac{x_jM'}{M}$ by $z^{\pm 1}$. Thus, every $\Gm$-orbit contains a unique point in the locus~\eqref{eq:yM'/M=1}. In other words, the locus~\eqref{eq:yM'/M=1} is isomorphic to the quotient $\Xcal(\Bice\pC)/\Gm$.  Since we assumed that $\alpha\pB_\ell = 1$, the quotient $\Xcal(\Bice\pC)/\Gm$ can be obtained by setting $x_\ell=1$.  Let $\Bice'$ be the $(n+m-2)\times(n-2)$ matrix obtained from $\Bice\pC$ by removing the $\ell$-th row. Let $\Qice'$ be the associated quiver, with $m$ frozen vertices corresponding to the rows $(\Vfro\sqcup\{j\})\setminus\{\ell\}$ and mutable part $Q-\{i,j\}$. Clearly, $\Qice'$ is still really full rank. We find that indeed the locus~\eqref{eq:yM'/M=1} is isomorphic to $\Xcal(\Qice')$.
%
%
  %
  %
\end{proof}

\begin{remark}
\Cref{prop:Louise_leaf} generalizes to the case of $i$ being either a source or a sink in $Q$.
\end{remark}
The following result will be later compared to the \FLY skein relation~\eqref{eq:HOMFLY_dfn}.
\begin{corollary}\label{cor:RQ_skein}
  Let $Q$ be really full rank, and let $i\in V(Q)$ be a \Louiseleaf in $Q$ connected to $j\in V(Q)$. Let $Q':=Q-\{i\}$ and $Q'':=Q-\{i,j\}$. Then
  \begin{equation}\label{eq:RQ_skein}
    \RQ(q)=(q-1)\RQX{Q'}{q} + q\RQX{Q''}{q}.
  \end{equation}
\end{corollary}
\begin{proof}
  Let $\Qice$ be really full rank with mutable quiver $Q$, and $\Xcal:=\XQice$. Let $\Ucal:=\{x_i\neq0\}$ and $\Zcal:=\{x_i=0\}$ be two subvarieties of $\Xcal$. Then $\Ucal$ is a really full rank cluster variety of type $Q'$ and $\Zcal$ is described in \cref{prop:Louise_leaf}. The result follows since
  \begin{equation*}%
    \#\Xcal(\F_q)=\#\Ucal(\F_q) + \#\Zcal(\F_q).    \qedhere
  \end{equation*}
\end{proof}
